% Langue 
\usepackage{polyglossia}
\setmainlanguage{french}

% Police de caractère
\usepackage{fontspec}   
\setmainfont[%
  Ligatures={TeX, Common, Historic, Rare},
  RawFeature = {+dlig,+hlig,+calt}
  ]{Libertinus Serif}
\usepackage{unicode-math}
\setmathfont{Erewhon-Math}
%\setmathfont{TeX Gyre Pagella Math}
%\setmathfont{Libertinus Math}

% Typographie
\usepackage[final]{microtype}
\allowdisplaybreaks
\sloppy
\setlength{\parindent}{0pt}
\setlength{\parskip}{3pt plus 0.5pt minus 0.5pt}

% Les couleurs
\usepackage{xcolor}
\definecolor{gris}{HTML}{808080}
\definecolor{Pantone7694C}{HTML}{01426A}

% Le package tikz
\usepackage{tikz}
\usetikzlibrary{%
    patterns, 
    decorations.fractals, 
    decorations.text, 
    angles, 
    quotes, 
    shapes.multipart, 
    arrows.meta, 
    calc, 
    tikzmark, 
    positioning, 
    shapes.geometric
}
\tikzset{
  gbox/.style={
    draw=gris, fill=gris!2,           
    line width=0.2mm, rounded corners=2mm, inner sep=2mm              
  }
}

% Le package tcolorbox
\usepackage[most]{tcolorbox}
\tcbset{
  general/.style={
    enhanced, breakable, drop small lifted shadow=gris,
    top=2mm, bottom=2mm, right=2mm, left=2mm, arc=2mm,
    boxrule=.2mm,
    colframe=gris, colback=gris!10, coltitle=gris, colbacktitle=gris!10
  }
}

% Des packages utiles
\usepackage{tabularray}
\usepackage{array}
\usepackage{graphicx}
\usepackage{xparse}
\usepackage{comment}

% Les symboles de délimitation
\newcommand*\intOO[2]{\left] #1 \, , \, #2 \right[}
\newcommand*\intFF[2]{\left[ #1 \, , \, #2 \right]}
\newcommand*\intOF[2]{\left] #1 \, , \, #2 \right]}
\newcommand*\intFO[2]{\left[ #1 \, , \, #2 \right[}
\newcommand*\intEN[2]{\left[\!\left[ #1 \, , \, #2 \right]\!\right]}
\newcommand*\spr[2]{\left< \, #1 \, \middle| \, #2 \, \right>}
\newcommand*\abs[1]{\left\lvert \, #1 \, \right\rvert}
\newcommand*\nor[1]{\left\lVert \, #1 \, \right\rVert}
\newcommand*\nop[1]{\left|\kern-0.15ex\left|\kern-0.15ex\left| \, #1 \, \right|\kern-0.15ex\right|\kern-0.15ex\right|}
\newcommand*\ket[1]{\left| \, #1 \, \right>}
\newcommand*\bra[1]{\left< \, #1 \, \right|}

% Le package enumitem
\usepackage{enumitem}
\setenumerate{topsep=1pt, itemsep=1pt, parsep=0pt, partopsep=0pt} 
\setitemize{topsep=1pt, itemsep=1pt, parsep=0pt, partopsep=0pt} 

% Le package hyperref
\usepackage{hyperref} 
\hypersetup{%
    colorlinks=true,
    linkcolor=couleurmatiere,
    urlcolor=couleurmatiere,
    filecolor=couleurmatiere,
    pdfauthor=Olivier Allard
    }